\section{Code--Data Identification and Data-Preserving Rewriting}
\label{sec:code_data}
\subsection{Problem Definition}
Let an input payload be a byte sequence $B=(b_0,\ldots,b_{n-1})$.
The identification task returns candidate data intervals
$[s_j,e_j)$ in payload-relative coordinates, optionally accompanied by
a kind such as an ASCII string, a UTF-16LE string, or another data object.
A candidate identifies a proposed role for the bytes. Its acceptance
for rewriting additionally depends on the available reference evidence
and the capabilities of the data-handling layer.

The original inline-data failure in Table~\ref{tab:broken-payload}
and the immediate-operand example in Table~\ref{tab:functional-payload}
motivate this extension. Both examples and their interpretation remain
in Section~\ref{sec:limitations}; the added layer explicitly represents
candidate data and the references affected by rewriting.

The task also includes uncertainty. A region outside the recovered graph
may contain data, unobserved code, or a mixture of the two. The pipeline
retains the origin of a candidate rather than collapsing all such
regions into the same evidential category. This permits separate
assessment of candidate generation, filtering, and downstream outcomes.

\subsection{Symbolic and Heuristic Identification}
The non-LLM path combines static graph information with bounded symbolic
exploration. A CFGFast pass supplies candidate code blocks and inferred
data objects. Symbolic exploration records instruction addresses visited
under the raw-payload model. An additional graph-recovery pass can use
newly observed entries, so that dynamically observed code is not lost
merely because it was absent from the first graph.

The symbolic initial state models uninitialized registers and memory
symbolically. It does not initialize a complete Windows process or
implement all external API behavior. Exploration is bounded by a time
budget, a state/block-step limit, a pending-state limit, and solver
budgets. Reports record the stop reason and unresolved states. Reaching
a resource limit, encountering an unsupported operation, and exhausting
the current frontier have different meanings and are reported
separately.

The implementation maintains evidence at byte granularity by expanding
observed instruction starts into their decoded spans. This avoids
treating the interior bytes of an observed instruction as unexecuted
data candidates. Candidate strings are then identified with printable
ASCII and UTF-16LE patterns, including their terminators. Overlap with
observed code prevents acceptance under this string heuristic.

Static graph recovery can itself classify a string as code. For a
string that overlaps candidate graph code, the implementation requires
additional structural corroboration rather than relying on printability
alone. Other proposed regions originate from graph-inferred data or
from bytes outside recovered blocks. The latter remain uncertain and
are distinguished from stronger data candidates.

The analysis returns intervals, their kinds, and diagnostic information;
it does not prove a complete partition of the payload into code and
data. In particular, a byte not observed during this run is not thereby
proven unreachable or safe to move. The experimental assessment records
both the proposed intervals and the limits of the run that produced them.

\subsection{LLM-Based Identification}
The alternative identification path presents the input through a hex
view, a linear disassembly, and candidate printable runs. These are
different views of the same bytes. The disassembly is explicitly
heuristic, and the printable runs are candidate evidence rather than
reference labels. The model returns structured data ranges, including
offset and length, instead of a rewritten binary.

The parser checks that intervals have positive lengths, remain within
the input bounds, and do not overlap. The engine then associates
candidate ranges with its existing reference information. Ranges that
cannot be accepted by the configured pipeline are reported or rejected.
A fresh engine instance is used when applying a proposal, avoiding
reuse of partially modified state after a rejected answer.

An invalid answer can trigger another request with feedback. The
reported outcome therefore distinguishes the first proposal, later
proposals, and the final accepted result. This distinction is necessary
to attribute an observation correctly: a useful final extraction may
reflect both the model's proposal and deterministic filtering.

The number of worked examples in this classification prompt is recorded
independently of the transformation experiment. A zero-shot/two-shot
comparison of assembly transformation does not imply that the same
comparison has been performed for identification.

\subsection{Shared Data and Reference Handling}
Both identification paths connect to deterministic data-handling
routines. These routines preserve the contents selected for extraction,
track the references whose interpretation depends on the data layout,
and finalize supported references after the code layout is known.
Identification answers which regions are candidates; reference handling
determines which proposed operations are supported by the implementation.

Different architectures require different addressing treatment. The
x64 and x86 implementations therefore share a conceptual contract
without being assumed identical at operand level. A candidate accepted
in one architecture need not be accepted in the other. The evaluation
keeps architecture and reference category associated with the outcome.

The pipeline also checks output integrity, including consistency of
the preserved data contents and the supported references. These checks
can prevent an invalid output from being accepted, but they cannot
establish that all references in an arbitrary binary have been recovered.
An unsupported case is consequently distinguished from a successfully
handled case and from an output that is built but fails execution.

Analysis caching is a practical part of the implementation. Candidate
data can be reused for an input identity, while references are recomputed
against the working representation. Reproducibility requires recording
whether an identification result was computed or retrieved and the
configuration associated with the cached entry. Timing comparisons
separate these cases.

\subsection{Classification Outputs and Residual Uncertainty}
Three outputs are retained for analysis: raw candidate intervals,
intervals accepted by deterministic checks, and intervals ultimately
processed by the data layer. Equal final binaries do not establish
equal raw classifications, because filtering can remove differences.
Conversely, a correct candidate can be rejected because its reference
pattern is unsupported. That rejection is an integration limitation,
not necessarily a classification error.

The additional evaluation in Section~\ref{sec:methodology} therefore separates
classification quality from downstream functionality. Successful
execution supplies evidence about the observed behavior after data
handling, while reference annotations supply evidence about the byte
roles. Neither observation can substitute for the other.
